% =============================================================================
% IV-E. Theoretical Analysis — density scaling, sparsity, residual, deployment
% Careful wording: density O(K/N); do NOT claim overall complexity is strictly linear
% =============================================================================

\subsection{Theoretical Analysis}
\label{sec:theory}

We do not claim global optimality of the joint objective.
We record structural properties that bind the method to the
scalability and deployment experiments in Sec.~\ref{sec:scale-deploy}.

\subsubsection{Adaptive Communication Optimization}
\label{sec:theory-obj}

With team return $\mathbb{E}[R]$ and soft activation mass $\mathbb{E}[C]$,
\begin{equation}
J(\theta)=\mathbb{E}[R]-\lambda\,\mathbb{E}[C],
\qquad
C=\sum_{i\neq j}g_{ij},
\label{eq:J}
\end{equation}
matches Eq.~\eqref{eq:joint-max}.
The decision variables include continuous gates $\{g_{ij}^{t}\}$ and the
budgeted topology $A_t^{\mathrm{AC}}$, co-optimized with $\pi_\theta$.
This differs from two-stage pruning:
$\theta^\star=\arg\max\mathbb{E}[R]$ followed by a separate reduction of $C$.

\subsubsection{Sparse Topology Complexity Bound}
\label{sec:theory-sparse}

\paragraph{Lemma~1 (Communication sparsification).}
Under a per-agent Top-$K$ budget with average degree $K$, the
\emph{attended} support has size $\mathcal{O}(NK)$ and message aggregation
costs $\mathcal{O}(NKd+Nd)$ (Sec.~\ref{sec:complexity}).
Relative to dense attention $\mathcal{O}(N^{2}d)$, the advantage widens when
$K$ remains bounded as $N$ grows.
Gate scoring on the radius-feasible set remains $\mathcal{O}(Nkd)$ and is
shared in spirit with geometric graph methods; the claimed gain concerns
\emph{activated} interaction density after gating/budgeting, not the elimination
of all candidate-edge construction.

\subsubsection{Communication Density Scaling}
\label{sec:theory-density}

\paragraph{Proposition~2 (Communication density scaling).}
\label{prop:density}
Define the normalized communication density
\begin{equation}
\eta_N
=
\frac{\sum_{i,j}g_{ij}}{N(N-1)}
=
\frac{C_N}{N(N-1)}.
\label{eq:eta}
\end{equation}
For a dense graph with $|E|=N(N-1)$ and $g_{ij}\equiv 1$,
$\eta_{\mathrm{dense}}=\mathcal{O}(1)$.
Under bounded effective activation degree $K\ll N$,
$C_N=\sum g_{ij}=\mathcal{O}(NK)$, hence
\begin{equation}
\eta_{\mathrm{AC}}
=
\mathcal{O}\!\left(\frac{K}{N}\right),
\qquad
\lim_{N\to\infty}\eta_{\mathrm{AC}}=0
\quad\text{($K$ bounded)}.
\label{eq:eta-limit}
\end{equation}
\begin{quote}
Under bounded neighborhood size $K$, the activated communication density
decreases inversely with swarm size.
\end{quote}
This theoretical trend is empirically verified in Fig.~\ref{fig:density-scale}
($N\in\{8,16,32\}$), where AC-DSGF maintains a near-constant sparse
activation density while DSGF exhibits substantially higher $\eta_N$.

\subsubsection{Residual Stability}
\label{sec:theory-residual}

\paragraph{Lemma~2 (Residual stability).}
If $\|\Delta(\Phi_i^{t})\|\le M$, then from Eq.~\eqref{eq:residual},
\begin{equation}
\bigl\|a_i^{t}-\pi_{\theta}(\cdot\mid o_i^{t})\bigr\|
\;\le\;
\beta(t)\,M.
\label{eq:residual-bound}
\end{equation}
As $\beta(t)\!\to\!0$, $a_i^{t}\!\to\!\pi_{\theta}(\cdot\mid o_i^{t})$.
Budgeted topology control cannot erase the base controller, explaining
graceful degradation under tight budgets (Sec.~\ref{sec:budget-res}).

\subsubsection{Deployment Approximation}
\label{sec:theory-deploy}

Training retains continuous $g_{ij}$ for differentiable exploration;
deployment requires discrete transmit decisions.
Per-agent Top-$K$ on the radius support,
\begin{equation}
A_{ij}^{\mathrm{deploy}}
=
\mathbf{1}\bigl[j\in\mathrm{Top}\text{-}K_i(g_{i\cdot})\bigr],
\label{eq:deploy}
\end{equation}
is an engineering discretization of the learned soft topology---not a
replacement of the training objective.
Sec.~\ref{sec:topk-deploy} shows that discrete links preserve task
effectiveness relative to soft inference.
